%%%%%%%%%%%%%%%%%%%%%%%%%%%%%%%%%%%%%%%%%%%%%%%%%%%%%%%%%%%%%%%%%%%%%%%%%%%%%%%%%%%%%%%%%%%%%%%%%%%%%%%%%%%%%%%%%%%%%%%%%%%%%%%%%%%%%%%%%%%%%%%%
%%%%%%%%%%%%%%%%%%%%%%%%%%%%%%%%%%%%%%%%%%%%%%%%%%%%%%%%%%%%%%%%%%%%%%%%%%%%%%%%%%%%%%%%%%%%%%%%%%%%%%%%%%%%%%%%%%%%%%%%%%%%%%%%%%%%%%%%%%%%%%%%
%%%%%%%%%%%%%%%%%%%%%%%%%%%%%%%%%%%%%%%%%%%%%%%%%%%%%%%%%%%%%%%%%%%%%%%%%%%%%%%%%%%%%%%%%%%%%%%%%%%%%%%%%%%%%%%%%%%%%%%%%%%%%%%%%%%%%%%%%%%%%%%%
%
%  Init Apollo-NG Beamer Slides
%


\documentclass{beamer}
\usepackage{tikz}
\usepackage{xcolor}
\usepackage{listings}
\usepackage{pxfonts}
\usepackage[english]{babel}
\usetikzlibrary{calc}
\usefonttheme{professionalfonts}
\usecolortheme[RGB={117,137,12}]{structure}

\definecolor{title}{RGB}{117,137,12}
\definecolor{block}{RGB}{185,182,174}
\definecolor{secinhead}{RGB}{227,220,201}
\definecolor{titlebg}{RGB}{51,51,51}
\setbeamercolor{secsubsec}{fg=secinhead}
\setbeamercolor{frametitle}{fg=secinhead}
\setbeamersize{text margin left=1.75em}

\setbeamertemplate{headline}
{
  \leavevmode%
  \hbox{%
  \begin{beamercolorbox}[wd=\paperwidth,ht=3ex,dp=2ex]{secsubsec}%
    \raggedright
    \hspace*{2em}%
    {\sffamily\fontsize{7}{9}\selectfont\color{secinhead}\thesection.~\insertsection\hfill\insertsubsection}%
    \hspace*{2em}%
  \end{beamercolorbox}%
  }%
}
\setbeamertemplate{frametitle}
{
  \leavevmode
  \hbox{%
  \begin{beamercolorbox}[wd=\paperwidth,ht=2.45ex,dp=0ex]{frametitle}%
    \raggedright\hspace*{1.3em}\Large\insertframetitle
  \end{beamercolorbox}
  }%
}

\setbeamercovered{transparent}

\newcommand\Fontvi{\fontsize{6}{7.2}\selectfont}

\setbeamertemplate{background canvas}{\includegraphics
	[width=\paperwidth,height=\paperheight]{background.png}}

\makeatletter
\definecolor{pbfc}{RGB}{117,137,12}% filling color for the progress bar
\definecolor{pbbg}{HTML}{282828}% background color for the progress bar
\def\progressbar@progressbar{} % the progress bar
\newcount\progressbar@tmpcounta% auxiliary counter
\newcount\progressbar@tmpcountb% auxiliary counter
\newdimen\progressbar@pbht %progressbar height
\newdimen\progressbar@pbwd %progressbar width
\newdimen\progressbar@tmpdim % auxiliary dimension

\progressbar@pbwd=70ex
\progressbar@pbht=1ex

% the progress bar
\def\progressbar@progressbar{%

    \progressbar@tmpcounta=\insertframenumber
    \progressbar@tmpcountb=\inserttotalframenumber
    \progressbar@tmpdim=\progressbar@pbwd
    \multiply\progressbar@tmpdim by \progressbar@tmpcounta
    \divide\progressbar@tmpdim by \progressbar@tmpcountb

  \begin{tikzpicture}[rounded corners=1pt,very thin]

    % Colors
    \colorlet{border}{black}

    \shade[top color=pbbg!100,bottom color=pbbg!80,top color=pbbg!100]
      (0pt, 0pt) rectangle ++ (\progressbar@pbwd, \progressbar@pbht);

      \shade[draw=pbfc,top color=pbfc!50,bottom color=pbfc!50,middle color=pbfc] %
        (0pt, 0pt) rectangle ++ (\progressbar@tmpdim, \progressbar@pbht);

      \draw[color=border]
      (0pt, 0pt) rectangle (\progressbar@pbwd, \progressbar@pbht);
  \end{tikzpicture}%
}


\setbeamertemplate{navigation symbols}{}
\setbeamercolor{normal text}{fg=white}

\addtobeamertemplate{frametitle}{\vskip-0.25ex}{}
\addtobeamertemplate{footline}{\vskip-0.5ex}{}

\addtobeamertemplate{footline}{}
{%
  \begin{beamercolorbox}[ht=3ex,leftskip=1.5ex,rightskip=1ex]{author in head/foot}
    \progressbar@progressbar  \hspace{1.5ex} \insertframenumber
    \vspace{1.5ex}
   \end{beamercolorbox}
}
\makeatother



%%%%%%%%%%%%%%%%%%%%%%%%%%%%%%%%%%%%%%%%%%%%%%%%%%%%%%%%%%%%%%%%%%%%%%%%%%%%%%%%%%%%%%%%%%%%%%%%%%%%%%%%%%%%%%%%%%%%%%%%%%%%%%%%%%%%%%%%%%%%%%%%
%%%%%%%%%%%%%%%%%%%%%%%%%%%%%%%%%%%%%%%%%%%%%%%%%%%%%%%%%%%%%%%%%%%%%%%%%%%%%%%%%%%%%%%%%%%%%%%%%%%%%%%%%%%%%%%%%%%%%%%%%%%%%%%%%%%%%%%%%%%%%%%%
%%%%%%%%%%%%%%%%%%%%%%%%%%%%%%%%%%%%%%%%%%%%%%%%%%%%%%%%%%%%%%%%%%%%%%%%%%%%%%%%%%%%%%%%%%%%%%%%%%%%%%%%%%%%%%%%%%%%%%%%%%%%%%%%%%%%%%%%%%%%%%%%
%
%  Begin Presentation
%



\title[DSpace]{\fontsize{22}{24}\selectfont DSpace}
\subtitle[A new way of handling geolocation based information ]{\fontsize{12}{14}\selectfont \textcolor{block}{A new way of handling geolocation based information }}
\author[iggy \& chrono]{\\[3ex]iggy \& chrono}

\institute[Apollo-NG]{
  \includegraphics[scale=0.5]{apollo-mission-badge.pdf}}
}
\date[June 2013]{June 13, 2013}


\begin{document}


%--- the titlepage frame -------------------------%
{
\setbeamertemplate{background canvas}{\includegraphics
	[width=\paperwidth,height=\paperheight]{background-title.png}}
\setbeamertemplate{headline}{}
\setbeamertemplate{footline}{}
\begin{frame}[plain]
  \vspace{1em}
  \titlepage
\end{frame}
}

\addtocounter{framenumber}{-1}

%--- the toc frame -------------------------%

{
\setbeamertemplate{headline}{}
\setbeamertemplate{footline}{}
\begin{frame}[plain]
\frametitle{\vspace{-2.12ex}Talk-Contents}
  \tableofcontents
\end{frame}
}

\addtocounter{framenumber}{-1}
\setcounter{page}{1}

%
%  Begin Content
%

\section{Section 1}
\begin{frame}{A sample slide}{with a subtitle}

A displayed ebola

\[
  \int_{-\infty}^\infty e^{-x^2} \, dx = \sqrt{\pi}
\]

An itemized list:

\begin{itemize}
  \item itemized item 1
  \item itemized item 2
  \item itemized item 3
\end{itemize}

\begin{theorem}
  In a right triangle, the square of hypotenuse equals
  the sum of squares of two other sides.
\end{theorem}

\end{frame}


\subsection{Uncovering Blockwise}
\setbeamertemplate{blocks}[rounded]
\setbeamercolor{block body}{bg=block}
\begin{frame}[fragile]
\frametitle{Inserting source code}
\Fontvi
\lstset{language=C++,
				numbers=left,
				numberstyle=\tiny,
				numbersep=-1.8ex,
                basicstyle=\color{black}\ttfamily,
                keywordstyle=\color{blue}\ttfamily,
                stringstyle=\color{red}\ttfamily,
                commentstyle=\color{green}\ttfamily,
                morecomment=[l][\color{magenta}]{\#}
}
\begin{block}{}
\begin{lstlisting}
    #include<stdio.h>
    #include<iostream>
    // A comment
    int main(void)
    {
    printf("Hello World\n");
    return 0;
    }
\end{lstlisting}
\end{block}
\end{frame}

\section{Sec 2}

\begin{frame}[The next Section...]

   \begin{columns}
        \column{.5\textwidth}
          A figure
        \column{.5\textwidth}
        \begin{overprint}
          \onslide<1>
          \begin{itemize}
              \item Item 1
              \item Item 2
          \end{itemize}
          \onslide<2>
          \begin{itemize}
              \item Another item 1
              \item Another item 2
              \item This list is longer
              \item Than the previous one
              \item was
          \end{itemize}
        \end{overprint}
      \end{columns}

\end{frame}
\end{document}
